\documentclass[10pt,a4paper]{amsart}
\usepackage[margin=19mm]{geometry}
\usepackage{amsmath,amssymb,mathtools}
\usepackage[scr=rsfs,cal=euler]{mathalfa}
\usepackage{xfrac}
\usepackage[unicode,hidelinks,bookmarks=false]{hyperref}
\newcommand{\ie}{i.e.\ }
\newcommand{\cf}{cf.\ }
\newcommand{\resp}{respectively\ }
\newcommand{\Subsection}{\subsection}
\providecommand{\nobreakdash}{\nobreak\hbox{-}}
\setlength{\emergencystretch}{2em}
\multlinegap0pt
\usepackage{bm}
\usepackage{bbm}
\usepackage{dsfont}
\usepackage{xspace}
\usepackage{cancel}
\usepackage{mathtools}
\usepackage{amsthm}
\makeatletter
\@ifundefined{thm}{\newtheorem{thm}{Thm}}{}
\@ifundefined{prop}{\newtheorem{prop}[thm]{Proposition}}{}
\makeatother
\newcommand{\M}{\mathbf{m}}
\newcommand{\N}{\mathbf{n}}
\title[Some more talents of the talented monoid of a higher-rank gr]{Some more talents of the talented monoid of a higher-rank graph}
\author{Roozbeh Hazrat, Huanhuan Li, Promit Mukherjee}
\date{Source: arXiv:2608.12798v1 (CC BY 4.0)}
\begin{document}
\maketitle

\noindent\emph{Source and licence.} Some more talents of the talented monoid of a higher-rank graph. Version arXiv:2608.12798v1 (CC BY 4.0), distributed under the Creative Commons Attribution 4.0 International licence, as recorded in the arXiv licence metadata for this exact version and shown on the abstract page https://arxiv.org/abs/2608.12798v1. Full rights notice: licenses/arxiv-2608.12798-CC-BY-4.0.txt.\par\smallskip

\noindent\textit{Editorial scope.} Counted unit: the Prop whose source label is \texttt{pro TM of product hrg} in the article (source0.tex lines 427--433, source label \texttt{pro TM of product hrg}), delivered together with the article's own complete original proof of it (source0.tex lines 434--479, one \texttt{proof} environment of 46 lines). No mathematics is written, repaired or re-derived: statement and proof are the article's own text line for line. Required context retained uncounted: none. Every definition, display and label of those lines is retained unchanged. Omitted from this package and not redistributed: 1--426, 480--1254. Nothing inside a retained excerpt is omitted or altered: every retained body is delivered exactly as the article's lines are written, so no omission receipt of any kind accompanies this package. Citations: the retained excerpts carry no citation command at all, so no bibliography entry of the article is transcribed and no reference is redistributed. Reference adaptation: none was needed and none was made; every reference inside the delivered unit resolves inside it, so no unresolved reference or citation is produced. Printed slips kept literal: no printed word, symbol, hypothesis or number of the counted statement or of its proof was altered. Layout and class: the article's own class is \texttt{amsart} with packages including fontenc, inputenc, amsmath, amssymb, amsfonts, amsthm, amscd, mathrsfs, bm, bbm, tikz, pgf, tikz-cd, tkz-graph; the wrapper is the lane's pinned \texttt{amsart} editorial wrapper at 10pt on A4, which restates the theorem-like environments the retained text needs and adds the article's own macro definitions that the retained text needs (\texttt{\textbackslash M}, \texttt{\textbackslash N}), with extra wrapper packages (bm, bbm, dsfont, xspace, cancel, mathtools). The delivered numbering is the unit's own. Assets: no retained span contains a figure environment, a graphics-inclusion command, a PDF-page inclusion, a table or a TikZ picture; no image file is copied into this package, the package declares no asset dependency and no image file is redistributed with it. Licence: version arXiv:2608.12798v1 of this article is distributed under the Creative Commons Attribution 4.0 International licence, as recorded in the arXiv licence metadata for this exact version and shown on the abstract page https://arxiv.org/abs/2608.12798v1. A full rights notice is delivered at licenses/arxiv-2608.12798-CC-BY-4.0.txt. Characters: every retained excerpt is pure ASCII, so the delivered engine, XeLaTeX with its default Latin Modern text font, reports no missing character.
\begin{prop}\label{pro TM of product hrg}
Let $(\Lambda_1,d_1)$ and $(\Lambda_2,d_2)$ be $k$- and $\ell$-graph respectively. Suppose both are row-finite and without sources. Then 

$(i)$ the tensor product $T_{\Lambda_1}\otimes T_{\Lambda_2}$ is a $\mathbb{Z}^{k+\ell}$-monoid, where \[^\mathbf{q}(v(\M)\otimes w(\N))=v(\M+\iota_1^{-1}(\pi_1(\mathbf{q})))\otimes w(\N+\iota_2^{-1}(\pi_2(\mathbf{q}))).\]

$(ii)$ $T_{\Lambda_1\times \Lambda_2}$ and $T_{\Lambda_1}\otimes T_{\Lambda_2}$ are isomorphic as $\mathbb{Z}^{k+\ell}$-monoids. 
\end{prop}

\begin{proof}
$(i)$ For each $\mathbf{q}\in \mathbb{Z}^{k+\ell}$, define a map 
\begin{align*}
\eta_\mathbf{q}:T_{\Lambda_1}\times T_{\Lambda_2}&\longrightarrow T_{\Lambda_1}\otimes T_{\Lambda_2}\\
\left(\displaystyle{\sum_{i=1}^{p}} v_i(\M_i),\displaystyle{\sum_{j=1}^{q}} w_j(\N_j)\right)&\longmapsto \displaystyle{\sum_{i=1}^{p}}\left(\displaystyle{\sum_{j=1}^{q}} v_i(\M_i+\iota_1^{-1}(\pi_1(\mathbf{q})))\otimes w_j(\N_j+\iota_2^{-1}(\pi_2(\mathbf{q})))\right). 
\end{align*}
One can then check that the map $\eta_\mathbf{q}$ is a well-defined bilinear map. As a consequence, it extends to a unique monoid homomorphism $\overline{\eta_\mathbf{q}}:T_{\Lambda_1}\otimes T_{\Lambda_2}\longrightarrow T_{\Lambda_1}\otimes T_{\Lambda_2}$. Obviously, $\overline{\eta_\mathbf{q}}$ is invertible with inverse homomorphism $\overline{\eta_{-\mathbf{q}}}$. Hence, $\mathbb{Z}^{k+\ell}$ acts on $T_{\Lambda_1}\otimes T_{\Lambda_2}$ by automorphisms, making it a $\mathbb{Z}^{k+\ell}$-monoid. It is evident that \[^\mathbf{q}(v(\M)\otimes w(\N))=\overline{\eta_\mathbf{q}}(v(\M)\otimes w(\N))=\eta_\mathbf{q}(v(\M),w(\N))=v(\M+\iota_1^{-1}(\pi_1(\mathbf{q})))\otimes w(\N+\iota_2^{-1}(\pi_2(\mathbf{q}))).\]

$(ii)$ Consider the following map: 
\begin{align*}
\phi:T_{\Lambda_1}\times T_{\Lambda_2}&\longrightarrow T_{\Lambda_1\times \Lambda_2}\\
\left(\displaystyle{\sum_{i=1}^{p}} v_i(\M_i),\displaystyle{\sum_{j=1}^{q}} w_j(\N_j)\right)&\longmapsto \displaystyle{\sum_{i=1}^{p}}\left(\displaystyle{\sum_{j=1}^{q}} (v_i,w_j)(\iota_1(\M_i)+\iota_2(\N_j))\right).
\end{align*}
Let $(v(\M),w(\N))\in T_{\Lambda_1}\times T_{\Lambda_2}$. In $T_{\Lambda_1}$, we have $v(\M)=\displaystyle{\sum_{\alpha\in v\Lambda_1^{\mathbf{q}}}} s(\alpha)(\M+\mathbf{q})$ for all $\mathbf{q}\in \mathbb{N}^k$; whereas in $T_{\Lambda_2}$, $w(\N)=\displaystyle{\sum_{\beta\in w\Lambda_2^{\mathbf{t}}}} s(\beta)(\N+\mathbf{t})$ for all $\mathbf{t}\in \mathbb{N}^\ell$. Note that 
\begin{align*}
&~(v,w)(\Lambda_1\times \Lambda_2)^{\iota_1(\mathbf{q})+\iota_2(\mathbf{t})}\\
=&~\{(\alpha,\beta)\in \Lambda_1\times \Lambda_2~|~r(\alpha)=v,r(\beta)=w,~\text{and}~\iota_1(d_1(\alpha))+\iota_2(d_2(\beta))=\iota_1(\mathbf{q})+\iota_2(\mathbf{t})\}\\
=&~\{\alpha\in \Lambda_1~|~r(\alpha)=v,\iota_1(d_1(\alpha))=\iota_1(\mathbf{q})\}\times \{\beta\in \Lambda_2~|~r(\beta)=w,\iota_2(d_2(\beta))=\iota_2(\mathbf{t})\}\\
=&~v\Lambda_1^{\mathbf{q}}\times w\Lambda_2^{\mathbf{t}}
\end{align*} since $\iota_1,\iota_2$ are injective. In view of this fact, we have 
\begin{align*}
&\phi\left(\displaystyle{\sum_{\alpha\in v\Lambda_1^{\mathbf{q}}}} s(\alpha)(\M+\mathbf{q}),\displaystyle{\sum_{\beta\in w\Lambda_2^{\mathbf{t}}}} s(\beta)(\N+\mathbf{t})\right)\\
=&\displaystyle{\sum_{\alpha\in v\Lambda_1^{\mathbf{q}},\beta\in w\Lambda_2^{\mathbf{t}}}} (s(\alpha),s(\beta))(\iota_1(\M+\mathbf{q})+\iota_2(\N+\mathbf{t}))\\
=&\displaystyle{\sum_{(\alpha,\beta)\in (v,w)(\Lambda_1\times \Lambda_2)^{\iota_1(\mathbf{q})+\iota_2(\mathbf{t})}}} s((\alpha,\beta))(\iota_1(\M)+\iota_2(\N)+\iota_1(\mathbf{q})+\iota_2(\mathbf{t}))\\
=&(v,w)(\iota_1(\M)+\iota_2(\N))=\phi(v(\M),w(\N)), 
\end{align*}
which shows that $\phi$ is well-defined. By definition, $\phi$ is bilinear. Hence, it extends to a unique monoid homomorphism $\Phi:T_{\Lambda_1}\otimes T_{\Lambda_2}\longrightarrow T_{\Lambda_1\times \Lambda_2}$. Now, we define a map in the other direction. Define $\Psi:T_{\Lambda_1\times \Lambda_2}\longrightarrow T_{\Lambda_1}\otimes T_{\Lambda_2}$ on generators as \[\Psi((v,w)(\mathbf{p})):=v(\iota_1^{-1}(\pi_1(\mathbf{p})))\otimes w(\iota_2^{-1}(\pi_2(\mathbf{p})))\] and then extend linearly to all of $T_{\Lambda_1\times \Lambda_2}$. Again, we check that $\Psi$ is well-defined. For any $\mathbf{q}\in \mathbb{N}^{k+\ell}$, using the same fact as used earlier, we have
\allowdisplaybreaks
{
\begin{align*}
&\Psi\left(\displaystyle{\sum_{(\lambda,\mu)\in (v,w)(\Lambda_1\times \Lambda_2)^{\mathbf{q}}}} s((\lambda,\mu))(\mathbf{p}+\mathbf{q})\right)\\
=&\displaystyle{\sum_{(\lambda,\mu)\in (v,w)(\Lambda_1\times \Lambda_2)^{\mathbf{q}}}} s(\lambda)(\iota_1^{-1}(\pi_1(\mathbf{p}+\mathbf{q})))\otimes s(\mu)(\iota_2^{-1}(\pi_2(\mathbf{p}+\mathbf{q})))\\
=&\displaystyle{\sum_{(\lambda,\mu)\in v\Lambda_1^{\iota_1^{-1}(\pi_1(\mathbf{q}))} \times w\Lambda_2^{\iota_2^{-1}(\pi_2(\mathbf{q}))}}} s(\lambda)(\iota_1^{-1}(\pi_1(\mathbf{p}))+\iota_1^{-1}(\pi_1(\mathbf{q})))\otimes s(\mu)(\iota_2^{-1}(\pi_2(\mathbf{p}))+\iota_2^{-1}(\pi_2(\mathbf{q})))\\
=&\left(\displaystyle{\sum_{\lambda\in v\Lambda_1^{\iota_1^{-1}(\pi_1(\mathbf{q}))}}} s(\lambda)(\iota_1^{-1}(\pi_1(\mathbf{p}))+\iota_1^{-1}(\pi_1(\mathbf{q})))\right)\otimes \left(\displaystyle{\sum_{\mu\in w\Lambda_2^{\iota_2^{-1}(\pi_2(\mathbf{q}))}}} s(\mu)(\iota_2^{-1}(\pi_2(\mathbf{p}))+\iota_2^{-1}(\pi_2(\mathbf{q})))\right)\\
=&~v(\iota_1^{-1}(\pi_1(\mathbf{p})))+w(\iota_2^{-1}(\pi_2(\mathbf{p})))=\Psi((v,w)(\mathbf{p})).
\end{align*}
}
Therefore, $\Psi$ is a well-defined monoid homomorphism. It is now easy to observe that $\Phi$ and $\Psi$ are mutually inverse homomorphisms. Thus, we have the desired isomorphism. Finally, for any generator $(v,w)(\mathbf{p})$ of $T_{\Lambda_1\times \Lambda_2}$, 
\begin{align*}
\Psi(^\mathbf{q}(v,w)(\mathbf{p}))&=\Psi((v,w)(\mathbf{p}+\mathbf{q}))\\
&=v(\iota_1^{-1}(\pi_1(\mathbf{p}+\mathbf{q})))\otimes w(\iota_2^{-1}(\pi_2(\mathbf{p}+\mathbf{q})))\\
&=^\mathbf{q}(v(\iota_1^{-1}(\pi_1(\mathbf{p})))\otimes w(\iota_2^{-1}(\pi_2(\mathbf{p})))\\
&=^\mathbf{q}\Psi((v,w)(\mathbf{p})).    
\end{align*}
Thus, the isomorphism $\Psi$ preserves the $\mathbb{Z}^{k+\ell}$-action and the result follows. 
\end{proof}

\end{document}
